% Ampliación propia de III; no modifica el núcleo común copiado de I.
% Residencia propuesta: antes de Retorno de hoja, incidencia areal--volumétrica.
% Procedencia y controles: VACANCIAS_PREFACTOR_PROCEDENCIA.md, en esta carpeta.
\section{Polarización discreta y modo trítico de las proyecciones aritméticas}
\label{iii:sec:vacancias-prefactor}

El refinamiento y la comparación entre las hojas de APP proporcionan dos
observaciones complementarias. La primera sigue la discrepancia acumulada
entre el número de transiciones y la capacidad adquirida; la segunda conserva
el modo del selector trítico en el emparejamiento de las fibras aditivas y
multiplicativas. La polarización discreta y el prefactor electrónico se
obtienen de estas operaciones respectivas. Su articulación con el vacío
requiere mantener tanto la información temporal como los espacios sobre los
que actúan los operadores.

\subsection{Discrepancia del refinamiento y convención de extremos}

Los cardinales del bloque ternario y de su representación decimal fijan
\[
 \rho=\log_{1000}729=\log_{10}9,\qquad
 \delta=1-\rho=\log_{10}(10/9).
\]
Se tiene $0<\delta<1$. Además, $\delta$ es irracional: una igualdad
$\delta=p/q$, con $q>0$, implicaría $10^{q-p}=9^q$, incompatible con
la descomposición en factores primos. Este desajuste ya interviene en la
codificación del refinamiento. Aquí se explicita su lectura centrada,
conservando el origen de fase y la convención de extremos.

\begin{definition}[Codificaciones inferior y superior]
Para $\theta\in\mathbb R$ y $n\in\mathbb Z$, se definen
\begin{align}
 z_n^{\downarrow}(\theta)
 &=\lfloor(n+1)\delta+\theta\rfloor-\lfloor n\delta+\theta\rfloor,
 \label{iii:vp:lower}\\
 z_n^{\uparrow}(\theta)
 &=\lceil(n+1)\delta+\theta\rceil-\lceil n\delta+\theta\rceil.
 \label{iii:vp:upper}
\end{align}
Ambas toman valores en $\{0,1\}$. La flecha distingue exclusivamente la
convención de extremos de los intervalos; la orientación trítica conserva
su definición anterior.
\end{definition}

La codificación inferior coincide con la utilizada en el desarrollo
del refinamiento. La superior permite recuperar la formulación de la
polarización con función techo. Las dos expresiones se relacionan por
un término de frontera exactamente determinado.

\begin{proposition}[Cambio de convención y término de frontera]
Sea $b_n(\theta)=\mathbf1_{\mathbb Z}(n\delta+\theta)$. Entonces
\begin{equation}
 z_n^{\uparrow}(\theta)-z_n^{\downarrow}(\theta)
 =b_n(\theta)-b_{n+1}(\theta).
 \label{iii:vp:endpoint}
\end{equation}
En particular, para $N\geq0$,
\[
 \sum_{n=0}^{N-1}(z_n^{\uparrow}-z_n^{\downarrow})
 =b_0-b_N.
\]
\end{proposition}
\begin{proof}
Para todo real $x$, $\lceil x\rceil=\lfloor x\rfloor+1-
\mathbf1_{\mathbb Z}(x)$. Aplicar esta identidad a los dos extremos
de cada incremento da \eqref{iii:vp:endpoint}; su suma es telescópica.
La irracionalidad de $\delta$ implica que, fijada $\theta$, como máximo
un índice entero satisface $n\delta+\theta\in\mathbb Z$.
\end{proof}

Con $\theta=0$, el primer símbolo superior es $1$ y el inferior es $0$;
para $n\geq1$ coinciden. El término inicial pertenece al registro y se
conserva cuando se cambia la convención. Las densidades y las cotas de
discrepancia se transportan con ese término, sin modificar la actualización
del TPK ni reiniciar la historia en cada retorno de fase.

\subsection{Polarización acotada y balance temporal exacto}

\begin{definition}[Polarización discreta centrada]
La polarización de la codificación superior, respecto de la densidad
$\delta$, es
\begin{equation}
 P_N^{\uparrow}(\theta)
 =\sum_{n=0}^{N-1}\bigl(z_n^{\uparrow}(\theta)-\delta\bigr),
 \qquad P_0^{\uparrow}=0.
 \label{iii:vp:polarization}
\end{equation}
Su incremento temporal es
\begin{equation}
 J_N^{\uparrow}(\theta)
 =P_{N+1}^{\uparrow}(\theta)-P_N^{\uparrow}(\theta)
 =z_N^{\uparrow}(\theta)-\delta.
 \label{iii:vp:current}
\end{equation}
\end{definition}

El centrado separa el crecimiento uniforme de densidad $\delta$ de la
información de fase. Así, $P_N^{\uparrow}$ mide un desequilibrio acumulado,
no el número total de vacancias. El incremento $J_N^{\uparrow}$ conserva
el pulso local y el fondo respecto del cual se evalúa.

\begin{theorem}[Cota uniforme y continuidad discreta temporal]
Para $N\geq0$ y toda fase $\theta$,
\begin{equation}
 P_N^{\uparrow}(\theta)
 =\lceil N\delta+\theta\rceil-\lceil\theta\rceil-N\delta,
 \qquad |P_N^{\uparrow}(\theta)|<1.
 \label{iii:vp:bound}
\end{equation}
Para cualesquiera enteros $0\leq a<b$,
\begin{equation}
 \sum_{n=a}^{b-1}J_n^{\uparrow}(\theta)
 =P_b^{\uparrow}(\theta)-P_a^{\uparrow}(\theta),
 \qquad
 \left|\sum_{n=a}^{b-1}J_n^{\uparrow}(\theta)\right|<1.
 \label{iii:vp:continuity}
\end{equation}
La densidad de la sucesión de vacancias es $\delta$.
\end{theorem}
\begin{proof}
La suma de los incrementos de techo de \eqref{iii:vp:upper} produce la
primera identidad. Escribiendo $r(x)=\lceil x\rceil-x\in[0,1)$,
se obtiene $P_N^{\uparrow}=r(N\delta+\theta)-r(\theta)$; su valor
absoluto es menor que uno. Sobre el intervalo $[a,b)$ se obtiene
análogamente $r(b\delta+\theta)-r(a\delta+\theta)$, lo que prueba
la segunda cota, más precisa que la suma de dos cotas individuales.
Finalmente,
$N^{-1}\sum_{n=0}^{N-1}z_n^{\uparrow}
=\delta+N^{-1}P_N^{\uparrow}\to\delta$.
\end{proof}

La codificación inferior admite las mismas demostraciones sustituyendo
$r(x)$ por $\lfloor x\rfloor-x\in(-1,0]$. En particular,
$P_N^{\uparrow}-P_N^{\downarrow}=b_0-b_N$. La discrepancia acotada
utilizada anteriormente en las lecturas de Dirichlet y la polarización
presente conservan, por tanto, una relación explícita de origen y extremos.
La aperiodicidad tampoco introduce una deriva secular de este desequilibrio:
la historia puede prolongarse indefinidamente mientras la discrepancia
centrada permanece uniformemente acotada.

\subsection{Realización en las rectas de carga y tiempo}

Para realizar dimensionalmente el balance se especifican una recta de carga
$\mathcal L_Q$ con sección no nula $u_Q$ y una sucesión creciente de
instantes $t_N$ en una recta temporal orientada. Con
$\Delta t_N=t_{N+1}-t_N>0$, la carga centrada y la corriente media de
cada intervalo quedan definidas por
\begin{equation}
 Q_N=P_N^{\uparrow}u_Q,\qquad
 I_N=\frac{Q_{N+1}-Q_N}{\Delta t_N}
     =\frac{u_Q}{\Delta t_N}\bigl(z_N^{\uparrow}-\delta\bigr).
 \label{iii:vp:dimensional}
\end{equation}
Así, $Q_N\in\mathcal L_Q$ e
$I_N\in\mathcal L_Q\otimes\mathcal L_T^{-1}$. La sección $u_Q$ y
los intervalos temporales son datos de esta realización; la construcción
aritmética precedente ya ha determinado $P_N^{\uparrow}$ y
$J_N^{\uparrow}$.

\begin{proposition}[Conservación del balance bajo realización dimensional]
Para $a<b$,
\begin{equation}
 Q_b-Q_a=\sum_{n=a}^{b-1}\Delta t_n I_n,
 \qquad
 \left|\frac{Q_b-Q_a}{u_Q}\right|<1.
 \label{iii:vp:dimensional-balance}
\end{equation}
Si todos los intervalos tienen longitud $t_0$, las dos corrientes posibles
son $-\delta u_Q/t_0$ y $(1-\delta)u_Q/t_0$.
\end{proposition}
\begin{proof}
Multiplicar cada identidad de \eqref{iii:vp:dimensional} por
$\Delta t_n$ y sumar telescopa las cargas. La cota es
\eqref{iii:vp:continuity}; los dos valores locales se obtienen de
$z_n^{\uparrow}\in\{0,1\}$.
\end{proof}

La continuidad demostrada es temporal y está referida a la carga centrada.
Una realización espacial incorpora además la localización de las cargas,
la incidencia entre celdas y la identificación de los flujos de frontera.
En una descripción constitutiva por canales, el acoplamiento se especifica
mediante una aplicación de excitación
$\iota:\mathcal L_Q\to\mathcal H_{\rm ch}\otimes\mathcal L_Q$
y una aplicación de observación
$\ell:\mathcal H_{\rm ch}\otimes\mathcal L_Q\to\mathcal L_Q$,
donde $\mathcal H_{\rm ch}$ es una realización real del espacio de canales.
Si ambas son
lineales e independientes del índice temporal, la respuesta de un operador
$\mathcal V$ fijo se transporta por
\[
 Q_N^{\rm resp}=\ell(\mathcal V\otimes I_{\mathcal L_Q})\iota(Q_N),
 \qquad
 I_N^{\rm resp}=\frac{Q_{N+1}^{\rm resp}-Q_N^{\rm resp}}{\Delta t_N}.
\]
La linealidad conserva el balance telescópico. Esta composición hace
explícitos los datos de una identificación constitutiva: la secuencia
escalar determina el pulso, mientras $\iota$, $\mathcal V$ y $\ell$
determinan su excitación, transporte y observación por canales. Las
identidades \eqref{iii:vp:polarization}--\eqref{iii:vp:dimensional-balance}
establecen la primera construcción con independencia de esa elección de
realización. El operador del vacío se construye en la
sección~\ref{iii:sec:constitutiva} sobre sus canales y proyectores propios.

\subsection{Emparejamiento de las fibras aditivas y multiplicativas}

La segunda observación procede de la realización tridimensional de APP,
$\mathcal A_3=D_9^3$, con medida uniforme y
$D_9=\{1,\ldots,9\}$. Se conservan los observables
\[
 \Sigma(x)=\rho_9(x_1+x_2+x_3),\qquad
 \Pi(x)=\rho_9(x_1x_2x_3),
\]
donde $\rho_9(n)=1+[(n-1)\bmod9]$ para $n>0$. En
$\ell^2(\mathcal A_3)$, sean $P_{\Sigma,3}$ y $P_{\Pi,3}$ las
esperanzas condicionales sobre sus fibras. Por ejemplo,
\[
 (P_{\Sigma,3}f)(x)=
 \frac1{|\Sigma^{-1}(\Sigma(x))|}
 \sum_{y:\,\Sigma(y)=\Sigma(x)}f(y).
\]
Son proyecciones ortogonales definidas sobre el mismo soporte. La matriz
de incidencias $N_{ab}=|\Sigma^{-1}(a)\cap\Pi^{-1}(b)|$ resulta de
enumerar los $729$ triples, sumando y multiplicando sus coordenadas antes
de aplicar $\rho_9$:
\begin{equation}
 N=9\begin{pmatrix}
 0&1&1&0&1&1&0&1&4\\
 1&0&1&1&0&1&1&0&4\\
 1&0&2&0&1&2&0&0&3\\
 0&1&1&0&1&1&0&1&4\\
 1&0&1&1&0&1&1&0&4\\
 0&0&2&1&0&2&0&1&3\\
 0&1&1&0&1&1&0&1&4\\
 1&0&1&1&0&1&1&0&4\\
 0&1&2&0&0&2&1&0&3
 \end{pmatrix}.
 \label{iii:vp:joint-matrix}
\end{equation}
Las filas suman $81$ y las sumas de columna son
\[
 (m_1,\ldots,m_9)=(36,36,108,36,36,108,36,36,297).
\]
Por ello, el emparejamiento de las indicatrices normalizadas de las fibras
es $C_{ab}=N_{ab}/\sqrt{81m_b}$. La normalización recoge las
multiplicidades producidas por APP, no una elección de valores singulares.

\begin{proposition}[Modo singular del selector trítico]
En la coordenatización canónica, el vector del selector es
$m_{\rm ph}=(1,-1,0)^3$. Se cumple
\begin{equation}
 Cm_{\rm ph}=C^{\mathsf T}m_{\rm ph}=-\tfrac12m_{\rm ph}.
 \label{iii:vp:trit-mode}
\end{equation}
Así, el valor singular correspondiente es $s=1/2$.
\end{proposition}
\begin{proof}
En las seis columnas donde $m_{\rm ph}$ es no nulo, $m_b=36$ y
$C_{ab}=N_{ab}/54$. Multiplicar las filas de
\eqref{iii:vp:joint-matrix} por $(1,-1,0)^3$ produce
$-m_{\rm ph}/2$. Para el producto transpuesto, las sumas firmadas
de las columnas $3,6,9$ son cero; las otras seis dan los mismos
valores $-m_{{\rm ph},b}/2$. Esto prueba ambas identidades y, por
composición, $CC^{\mathsf T}m_{\rm ph}=m_{\rm ph}/4$.
\end{proof}

\subsection{Norma del conmutador y prefactor electrónico}

\begin{theorem}[Norma exacta de incompatibilidad aritmética]
Las dos proyecciones satisfacen
\begin{equation}
 \bigl\|[P_{\Sigma,3},P_{\Pi,3}]\bigr\|=\frac{\sqrt3}{4}.
 \label{iii:vp:prefactor}
\end{equation}
El máximo se alcanza en el plano principal correspondiente al modo
trítico de \eqref{iii:vp:trit-mode}.
\end{theorem}
\begin{proof}
La matriz racional $M=CC^{\mathsf T}$ tiene entradas
$M_{ac}=\sum_bN_{ab}N_{cb}/(81m_b)$. Su espectro completo se
comprueba mediante los siguientes subespacios independientes: los
vectores $(1,\ldots,1)$, $(1,-1,0)^3$ y $(1,1,-2)^3$ tienen
autovalores $1$, $1/4$ y $7/132$, respectivamente; el subespacio de
vectores soportados en $\{3,6,9\}$ cuya suma es cero tiene dimensión
$2$ y autovalor $1/18$; los vectores de suma cero soportados en
$\{1,4,7\}$ o en $\{2,5,8\}$ forman un núcleo de dimensión $4$.
La sustitución en la matriz verifica estas acciones y las dimensiones
suman $9$. En consecuencia,
\[
 \operatorname{spec}M=
 \{1,1/4,7/132,1/18,1/18,0,0,0,0\}.
\]
Para dos proyecciones ortogonales, un valor singular $s\in(0,1)$ del
emparejamiento determina un plano principal en el que sus matrices son
\[
 P=\begin{pmatrix}1&0\\0&0\end{pmatrix},\qquad
 Q=\begin{pmatrix}s^2&s\sqrt{1-s^2}\\
 s\sqrt{1-s^2}&1-s^2\end{pmatrix}.
\]
Esta forma se obtiene tomando un vector unitario $u$ de la primera
imagen con $PQu=s^2u$ y completándolo con
$(Qu-s^2u)/(s\sqrt{1-s^2})$. Su conmutador es
$s\sqrt{1-s^2}\bigl(\begin{smallmatrix}0&1\\-1&0\end{smallmatrix}\bigr)$.
Los sectores de valor singular $0$ o $1$ contribuyen con cero.
Por tanto, la norma total es el máximo de
$\sqrt{\lambda(1-\lambda)}$ sobre el espectro de $M$. Dicho máximo
es $\sqrt{3/16}$, alcanzado en $\lambda=1/4$.
\end{proof}

La norma es el prefactor que interviene en la composición electrónica.
Su obtención conserva el soporte tridimensional, las fibras y el modo
del selector; la evaluación escalar no sustituye esa estructura. En
particular, conocer $\sqrt3/4$ permite conocer la magnitud máxima de
la incompatibilidad, mientras los proyectores conservan las direcciones
en que actúa.

\clearpage
\subsection{Defecto cuadrático y respuesta de los sectores \(90/120\)}

El modo trítico determina simultáneamente
\begin{equation}
 1-s^2=\frac34=\frac{90}{120},\qquad s=\frac12.
 \label{iii:vp:defect-ratio}
\end{equation}
La primera expresión es el defecto cuadrático del plano principal de APP.
La última conserva las longitudes de los dos sectores del calendario.
La respuesta constitutiva utiliza esos sectores sobre otro objeto: un
canal contractivo $q\in(0,1)$. Su evaluación es
\[
 r(q)=\frac{1-q^{90}}{1-q^{120}}.
\]
La relación entre el cociente de sectores y esta respuesta puede
precisarse sin identificar sus operadores.

\begin{proposition}[Límite del cociente de respuesta]
Para $0<q<1$,
\begin{equation}
 \frac34<r(q)<1,\qquad
 \lim_{q\uparrow1}r(q)=\frac34=1-s^2.
 \label{iii:vp:response-limit}
\end{equation}
\end{proposition}
\begin{proof}
Con $u=q^{30}\in(0,1)$, la factorización de $1-u^3$ y
$1-u^4$ da
$r(q)=(1+u+u^2)/(1+u+u^2+u^3)$. El límite en $u=1$ es
$3/4$. Además,
$4(1+u+u^2)-3(1+u+u^2+u^3)
=(1-u)(3u^2+2u+1)>0$, y el denominador supera al numerador
en $u^3>0$. Se obtienen las dos desigualdades.
\end{proof}

La coincidencia exacta concierne al defecto del modo singular y al límite
de la respuesta de los sectores. En el interior contractivo, la respuesta
depende de $q$ y es estrictamente mayor que $3/4$. El desarrollo
constitutivo conserva sus dos canales $q_\pm$ y sus proyectores
$P_\pm$; los proyectores $P_{\Sigma,3}$ y $P_{\Pi,3}$ conservan,
por su parte, las fibras aritméticas que determinaron el prefactor.
Una equivalencia operatoria entre ambas realizaciones exige una aplicación
entre sus espacios que entrelace las acciones correspondientes.
Las ecuaciones \eqref{iii:vp:trit-mode}, \eqref{iii:vp:prefactor} y
\eqref{iii:vp:response-limit} precisan la información compartida y
mantienen visibles los datos de cada realización.
